\Askhsh{21185}{4}
\begin{ylh}{1}
    \item 9.2. Το Πυθαγόρειο θεώρημα
    \item 9.4. Γενίκευση του Πυθαγόρειου θεωρήματος.
\end{ylh}
Τρία ευθύγραμμα τμήματα $\alpha$, $\beta$ και $\gamma$ έχουν μήκη ανάλογα των αριθμών 5, 4 και 3 αντίστοιχα.
\begin{alist}
\item Να αποδείξετε ότι τα ευθύγραμμα τμήματα αυτά μπορούν να σχηματίσουν ορθογώνιο τρίγωνο. \monades{8}
\item Αν τα ευθύγραμμα τμήματα $\alpha$, $\beta$ και $\gamma$ είναι σχεδιασμένα πάνω σε ένα χαρτί και αυτό το φωτοτυπήσουμε με μεγέθυνση $\lambda$\%, να αποδείξετε ότι και με τα νέα ευθύγραμμα τμήματα σχηματίζεται πάλι ορθογώνιο τρίγωνο. \monades{10}
\item Να εξετάστε αν μπορεί να σχηματιστεί τρίγωνο με πλευρές $10\alpha$, $8\beta$ και $6\gamma$. \monades{7}
\end{alist}